% !TEX root = ../hw01.tex
Put $I=(0,1)$. For each number in $I$, we choose the unique
decimal expansion that is not eventually all $9$s.
For example, we write one half as $0.5000\ldots$ instead of
$0.4999\ldots$. Each digit may be any of $0,1,\ldots,9$.

For $x=0.x_1x_2\ldots$ and $y=0.y_1y_2\ldots$, define
\[
  F(x,y)=0.x_1y_1x_2y_2\ldots.
\]
Figure~\ref{fig:hw01:decimal-interleaving} shows how the odd
positions record $x$ and the even positions record $y$.

\begin{figure*}[h]
  \centering
  % !TEX root = ../../problems/hw01.tex
\begin{tikzpicture}[x=1cm,y=1cm,>=Stealth,font=\small]
  \node[anchor=east] at (-0.75,1.25) {$x=0.$};
  \node[anchor=east] at (-0.75,0) {$F(x,y)=0.$};
  \node[anchor=east] at (-0.75,-1.25) {$y=0.$};
  \foreach \index in {1,2,3,4} {
    \pgfmathsetmacro{\oddpos}{3.0*(\index-1)}
    \pgfmathsetmacro{\evenpos}{\oddpos+1.5}
    \node[text=black] (x\index) at (\oddpos,1.25) {$x_{\index}$};
    \node[text=black] (y\index) at (\evenpos,-1.25) {$y_{\index}$};
    \node[fill=black!6,minimum width=0.65cm,minimum height=0.5cm]
      (odd\index) at (\oddpos,0) {$x_{\index}$};
    \node[fill=black!12,minimum width=0.65cm,minimum height=0.5cm]
      (even\index) at (\evenpos,0) {$y_{\index}$};
    \draw[->,black!65] (x\index.south) -- (odd\index.north);
    \draw[->,black!65] (y\index.north) -- (even\index.south);
  }
  \foreach \height in {-1.25,0,1.25}
    \node at (12.0,\height) {$\cdots$};
\end{tikzpicture}

  \caption{Interleaving keeps both digit sequences. Read the odd
    positions to recover $x$, and the even positions to recover $y$.}
  \label{fig:hw01:decimal-interleaving}
\end{figure*}

\clearpage

We first check that this construction gives a number in $I$.
The output digits cannot be eventually all $9$s. Otherwise,
both input expansions would also be eventually all $9$s,
contradicting our choice of decimal expansions.

Since $x>0$, at least one of its digits is nonzero. This digit
occurs in an odd position of the output, so $F(x,y)>0$.
The output is less than $1$ because its digits are not all $9$s.
Therefore $F(x,y)\in I$, and its displayed decimal expansion
follows our convention.

If $F(x,y)=F(u,v)$, uniqueness of the digits gives $x_j=u_j$
and $y_j=v_j$ for every $j$. Hence $x=u$ and $y=v$, so $F$
is one-to-one.

However, $F$ is not onto, despite the claim in the original hint.
To see this, consider $z=0.101010\ldots\in I$.
If $F(x,y)=z$, the even digits would force $y=0.000\ldots=0$.
This is impossible because $0\notin I$, so $z$ has no preimage.

We therefore use an injection in the reverse direction as well.
Define a map $G$ from $I$ to $I\times I$ by
\[
  G(t)=\left(t,\frac12\right).
\]
Both coordinates lie in $I$, so the map is well-defined.
If $G(t)=G(s)$, equality of the first coordinates gives $t=s$.
Hence $G$ is one-to-one.

The Cantor--Schr\"oder--Bernstein theorem says that injections
in both directions imply the existence of a bijection.
Applying it to $F$ and $G$ gives $|I\times I|=|I|$.
By Exercise~\ref{ex:hw01:interval-to-reals}, $|I|=|\R|=\aleph$.
Therefore
\[
  |(0,1)\times(0,1)|=|(0,1)|=\aleph.\qedhere
\]
